\documentclass[12pt]{article}

\usepackage{amsmath, amssymb, amsthm}
\usepackage{geometry}
\usepackage{tikz}
\usepackage{hyperref}
\usepackage{enumitem}
\usepackage{microtype}

\geometry{a4paper, margin=2.5cm}

\hypersetup{
    colorlinks=true,
    linkcolor=blue,
    urlcolor=blue
}

%% Theorem environments
\theoremstyle{definition}
\newtheorem{definition}{Definition}[section]
\newtheorem{example}[definition]{Example}

\theoremstyle{plain}
\newtheorem{theorem}{Theorem}[section]

\theoremstyle{remark}
\newtheorem{remark}{Remark}[section]

%% Shortcuts
\newcommand{\N}{\mathbb{N}}
\newcommand{\R}{\mathbb{R}}
\newcommand{\iso}{\sim}
\newcommand{\notiso}{\not\sim}

% ----------------------------------------------------------------
\title{\textbf{The Set of Real Numbers is Uncountable}}
\author{0xHaru}
\date{\today}
% ----------------------------------------------------------------

\begin{document}
\maketitle

% ================================================================
\section{Cardinality and Isomorphism}
% ================================================================

\begin{definition}[Isomorphism]
Two sets $A$ and $B$ are \textbf{isomorphic} (or \textbf{equinumerous}) if and only if
there exists a bijection between them. In that case we write $A \sim B$.
\end{definition}

\begin{definition}[Same Cardinality]
Given two sets $A$ and $B$, if $A \sim B$ we say that $A$ and $B$ are
\textbf{equinumerous} or, equivalently, that they have the \textbf{same cardinality}.
\end{definition}

\begin{definition}[Finite Cardinality]
For every $n \in \N$ define the set
\[
    J_n \;=\;
    \begin{cases}
        \emptyset       & \text{if } n = 0,\\
        \{1,\ldots, n\} & \text{if } n > 0.
    \end{cases}
\]
A set $A$ has \textbf{finite cardinality} if and only if $\exists\, n\in\N$ such that
$A \sim J_n$.
\end{definition}

\begin{definition}[Infinite Cardinality]
A set $A$ has \textbf{infinite cardinality} if and only if $\nexists\, n\in\N$ such
that $A \sim J_n$ (i.e.\ $A$ does \emph{not} have finite cardinality).
\end{definition}

% ================================================================
\section{Countable and Uncountable Sets}
% ================================================================

\begin{definition}[Countable Set]
An infinite set $A$ is \textbf{countable} if and only if $A \sim \N$.
\end{definition}

\begin{example}
The following infinite sets are countable.
\begin{itemize}[nosep]
    \item $\N^{+} = \N \setminus \{0\}$: a bijection $f\colon \N \to \N^{+}$ is given by
          $f(n)=\mathrm{succ}(n)$.
    \item The set of even numbers $\mathbb{E}$: a bijection $f\colon \N \to \mathbb{E}$ is given by
          $f(n)=2n$.
\end{itemize}
\end{example}

\begin{definition}[Uncountable Set]
An infinite set $A$ is \textbf{uncountable} if and only if $A \not\sim \N$.
\end{definition}

% ================================================================
\section{\texorpdfstring{$\R$}{R} is Uncountable}
% ================================================================

\begin{theorem}
$\R$ is not countable, i.e.\ $\R \not\sim \N$.
\end{theorem}

\begin{proof}
The proof proceeds in three steps.

%------------------------------------------------------------------
\medskip
\noindent\textbf{Step~1.}\quad $\R \sim (a,b)$ for all $a,b\in\R$ with $a<b$.

\smallskip
We show geometrically that there exists a bijection between any open interval
$(a,b)$ and $\R$. Place the segment $[a,b]$ horizontally above the real line.
Draw the lower semicircle with diameter $[a,b]$ and centre $c$. Given any point
$x\in\R$, draw the line from $x$ to $c$; it meets the semicircle at a unique point,
from which a perpendicular to $[a,b]$ yields a unique point in $(a,b)$. This defines
a bijection $\R\to(a,b)$; choosing $a=0$ and $b=1$ gives $\R\sim(0,1)$.

\begin{center}
\begin{tikzpicture}[scale=1.2]
    % Real line
    \draw[->] (-3.5,0) -- (3.5,0) node[right] {$\R$};
    \fill (0,0) circle (1.5pt) node[below] {$0$};

    % Diameter [a,b] at height 3, radius 1.8
    \draw (-1.8,3) -- (1.8,3);
    \fill (-1.8,3) circle (1.5pt) node[above left]  {$a$};
    \fill ( 1.8,3) circle (1.5pt) node[above right] {$b$};
    \fill (   0,3) circle (1.5pt) node[above]        {$c$};

    % Lower semicircle: center (0,3), radius 1.8
    \draw (-1.8,3) arc (180:360:1.8);

    % Example point p = (-2.8, 0) on R (to the left)
    \fill (-2.8,0) circle (1.5pt) node[below] {$x$};

    % Dotted line 1: from p=(-2.8,0) to c=(0,3)
    % Intersection with semicircle: q ~ (-1.228, 1.684)
    \draw[dotted, thick] (-2.8,0) -- (0,3);

    % Intersection point q on semicircle
    \fill (-1.228,1.684) circle (1.5pt);

    % Dotted line 2: from q perpendicular (vertical) up to [a,b]
    \draw[dotted, thick] (-1.228,1.684) -- (-1.228,3);

    % Corresponding point on [a,b]
    \fill (-1.228,3) circle (1.5pt);
\end{tikzpicture}
\end{center}

\smallskip
\noindent\textit{How the bijection works in practice.}
\begin{itemize}[nosep]
    \item \textbf{From $\R$ to $(a,b)$:} pick any point $x$ on the real line.
          Draw the straight line from $x$ to the centre $c$; it crosses the
          semicircle at exactly one point. From that crossing point, drop a
          vertical line up to the segment $[a,b]$: the point where it lands is
          the image of $x$ in $(a,b)$.

    \item \textbf{From $(a,b)$ to $\R$:} pick any point $p$ in $(a,b)$.
          Drop a vertical line from $p$ down to the semicircle; it meets the
          arc at exactly one point. Draw the straight line from that arc point
          through $c$ and extend it until it hits the real line: that
          intersection is the image of $p$ in $\R$.
\end{itemize}

%------------------------------------------------------------------
\medskip
\noindent\textbf{Step~2.}\quad $\N \not\sim (0,1)$\quad(\emph{by diagonalisation}).

\smallskip
Suppose, for contradiction, that $(0,1)\sim\N$. Under this hypothesis,
$(0,1)$ can be listed exhaustively.  Consider then \emph{any} exhaustive list of
elements of $(0,1)$:
\[
\begin{array}{lllllll}
  0, & \underline{a_{00}} & a_{01} & a_{02} & a_{03} & \cdots \\[2pt]
  0, & a_{10} & \underline{a_{11}} & a_{12} & a_{13} & \cdots \\[2pt]
  0, & a_{20} & a_{21} & \underline{a_{22}} & a_{23} & \cdots \\[2pt]
     & \vdots
\end{array}
\]
where $a_{ij}\in\{0,1,\ldots,9\}$ is the $j$-th decimal digit of the $i$-th
element in the list.

Construct the real number $0{,}\,c_0\,c_1\,c_2\,\ldots$ by setting
\[
    c_i \;=\;
    \begin{cases}
        2 & \text{if } a_{ii} \neq 2,\\
        3 & \text{if } a_{ii} = 2.
    \end{cases}
\]

However, $0{,}\,c_0\,c_1\,c_2\,\ldots$ does \emph{not} appear in the list.
Indeed, for every $i\geq 0$:
\begin{align*}
    a_{00} \neq c_0 \quad&\Longrightarrow\quad
    0{,}\,c_0c_1c_2\ldots \;\neq\; 0{,}\,a_{00}\,a_{01}\,a_{02}\,\ldots\\
    a_{11} \neq c_1 \quad&\Longrightarrow\quad
    0{,}\,c_0c_1c_2\ldots \;\neq\; 0{,}\,a_{10}\,a_{11}\,a_{12}\,\ldots\\
    a_{22} \neq c_2 \quad&\Longrightarrow\quad
    0{,}\,c_0c_1c_2\ldots \;\neq\; 0{,}\,a_{20}\,a_{21}\,a_{22}\,\ldots\\
    &\;\vdots
\end{align*}
The $i$-th digit of $0{,}\,c_0c_1c_2\ldots$ always differs from the $i$-th digit of
the $i$-th element in the list, so the constructed number differs from \emph{every}
element in the list.

Therefore the list is \textbf{not exhaustive} --- a contradiction. $\bot$

%------------------------------------------------------------------
\medskip
\noindent\textbf{Step~3.}\quad $\R \not\sim \N$.

\smallskip
By Step~1, $\R \sim (0,1)$; by Step~2, $(0,1)\not\sim\N$. By the transitivity
of $\sim$ it follows that
\[
    \R \;\not\sim\; \N,
\]
i.e.\ $\R$ is uncountable.
\end{proof}

\end{document}
